% Канонический TikZ-рисунок варианта 06. Править только здесь.
\newcommand{\figParallelogramTwoBisectors}{%
\begin{tikzpicture}[scale=0.78,line cap=round,line join=round,every node/.style={font=\small}]
  \coordinate (A) at (0,0); \coordinate (B) at (6,0);
  \coordinate (D) at (1,2.828); \coordinate (C) at (7,2.828); \coordinate (E) at (3,0);
  \draw[line width=0.9pt] (A)--(B)--(C)--(D)--cycle;
  \draw[line width=0.9pt] (D)--(E)--(C);
  \pic[draw,line width=0.65pt,angle radius=4.8mm] {angle=A--D--E};
  \pic[draw,line width=0.65pt,angle radius=4.8mm] {angle=E--D--C};
  \pic[draw,line width=0.65pt,angle radius=4.8mm] {angle=D--C--E};
  \pic[draw,line width=0.65pt,angle radius=4.8mm] {angle=E--C--B};
  \node[below left] at (A) {$A$}; \node[below right] at (B) {$B$};
  \node[above right] at (C) {$C$}; \node[above left] at (D) {$D$}; \node[below] at (E) {$E$};
\end{tikzpicture}}
